% =====================================================================
% APPENDIX: METHOD DEFINITIONS
% Moved on 2026-11-05 from the former main-text section "Refusal-Removal
% Methods" (sections/03_related_work.tex, draft v0 2026-11-04) when that
% section was rewritten as a house-style Related Work section
% (\label{sec:related}). Content is kept verbatim apart from cross-reference
% wording; every original \label is kept (sec:methods, sec:m-*).
% Table tab:methods stays in the Related Work section.
% Must be \input after \appendix (e.g. from sections/08_appendix.tex).
% =====================================================================
\section{Method Definitions}
\label{app:method-defs}
\label{sec:methods}
\label{app:methods}% former template placeholder label from 08_appendix.tex (never referenced; kept)

This appendix gives the notation, the update rule and the rationale of each
refusal-removal method reviewed in Section~\ref{sec:related}, together with
the strengths and weaknesses of each and the way it is instantiated in this
study.

\paragraph{Notation.}
A decoder-only transformer maintains a residual stream
$\mathbf{x}^{(\ell)} \in \mathbb{R}^{d}$ at each layer $\ell$. Each layer reads
from the stream and writes back to it through two matrices:
the attention output projection $W_O^{(\ell)}$ and the MLP down-projection
$W_{\mathrm{down}}^{(\ell)}$. Let $\mathcal{D}_{\mathrm{harm}}$ and
$\mathcal{D}_{\mathrm{safe}}$ be sets of harmful and harmless prompts, and let
$\boldsymbol{\mu}^{(\ell)}_{\mathrm{harm}}$ and
$\boldsymbol{\mu}^{(\ell)}_{\mathrm{safe}}$ be the mean residual activations at
the last prompt token over each set.

Methods are grouped by \emph{what they assume about refusal}: (i) it is a single
linear direction, (ii) it is a low-dimensional subspace or cone, (iii) it is
whatever transformation of each module's outputs separates harmful from
harmless inputs (no geometric assumption), or (iv) it is best located in a
different representation space altogether. Table~\ref{tab:methods}
(Section~\ref{sec:related}) summarises
the evaluated methods as instantiated in this study (\S\ref{sec:setup-methods}).

% ---------------------------------------------------------------------
\subsection{Single-direction methods}

\subsubsection{Classic abliteration (difference-in-means)}
\label{sec:m-dim}
\textbf{Mechanism.} Compute the refusal direction
$\hat{\mathbf{r}}^{(\ell)} =
\mathrm{normalize}(\boldsymbol{\mu}^{(\ell)}_{\mathrm{harm}} -
\boldsymbol{\mu}^{(\ell)}_{\mathrm{safe}})$, pick one layer's direction
$\hat{\mathbf{r}}$, and orthogonalise every matrix that writes to the residual
stream against it:
$W \leftarrow W - \lambda\,\hat{\mathbf{r}}\hat{\mathbf{r}}^{\top} W$
\cite{arditi2024refusal}. With $\lambda=1$ the model can no longer write any
component along $\hat{\mathbf{r}}$.

\textbf{Why it works.} Under the linear representation hypothesis
\cite{park2024linear}, high-level features correspond to directions.
Arditi et al.~\cite{arditi2024refusal} showed causally that adding $\hat{\mathbf{r}}$
induces refusal on harmless prompts and ablating it suppresses refusal on
harmful ones, across 13 chat models. Because the edit is a rank-one projection,
most of the weight matrix is untouched.

\textbf{Why it exists.} It turned an interpretability finding into a
practical, training-free procedure that needs only forward passes on a few
hundred prompts \cite{labonne2024abliteration,failspy_abliterator}.

\textbf{Pros.} Simple, fast (minutes), interpretable, rank-one, few
hyperparameters (layer, strength).
\textbf{Cons.} Assumes refusal is one-dimensional; the direction also carries
part of the harmless mean, so removing it damages unrelated behaviour (refusal
is better described as an affine function whose offset plain ablation ignores;
on a recurrent RWKV model plain ablation produced nonsense where affine concept
editing stayed coherent \cite{marshall2024refusal}); changing row norms shifts
activation scales; quality depends heavily on the chosen layer and strength,
which are usually set by hand (in an earlier cross-tool comparison, GSM8K
accuracy changed by between $+1.51$ and $-18.81$ percentage points
\cite{young2025comparative}). DECCP, ErisForge and
FailSpy's tool are implementations of this recipe with different layer-selection
heuristics. In this study diff-in-means runs inside Heretic with the same
automated search over direction, strength and layer profile as MPOA
(``Heretic-tuned DiM'', \S\ref{sec:setup-methods}).

\subsubsection{Projected and magnitude-preserving abliteration (MPOA)}
\label{sec:m-mpoa}
\textbf{Mechanism.} Two refinements of \S\ref{sec:m-dim}.
(1) \emph{Projection}: remove from $\hat{\mathbf{r}}$ its component along the
normalised harmless mean, so that only the part of the refusal direction
orthogonal to ordinary activations is ablated \cite{lai2024projected}.
(2) \emph{Magnitude preservation}: row-normalise $W$, apply the ablation, then
restore each row's original norm \cite{lai2025mpoa}.
Heretic adds a search over a per-layer weight kernel (peak strength, peak
position, minimum strength, falloff distance) for attention and MLP separately,
global versus per-layer directions, and a fractional direction index
interpolated between layers, all tuned by multivariate TPE
\cite{akiba2019optuna,heretic2025}. The edit is stored as a low-rank LoRA
adapter \cite{hu2022lora}.

\textbf{Why it works.} Projection reduces collateral damage: the part of
$\hat{\mathbf{r}}$ shared with harmless activations is exactly what is
needed for normal behaviour. Norm preservation keeps downstream layer norms and
activation magnitudes in the range the model was trained on, so edits
change \emph{direction} rather than \emph{scale}.

\textbf{Why it exists.} Plain abliteration measurably degraded capability
(for example GSM8K drops of several points in an earlier comparison
\cite{young2025comparative}); these refinements recover much of that loss.

\textbf{Pros.} Better capability retention at the same refusal rate;
automated hyperparameter search removes manual tuning.
\textbf{Cons.} Still assumes a single direction per layer; norm preservation is
only approximate when stored as a low-rank adapter; the search adds many
evaluation passes (one per trial).

\subsubsection{Null-space-constrained abliteration (jwest33)}
\label{sec:m-null}
\textbf{Mechanism.} Compute a projected, norm-rescaled diff-in-means edit
$\Delta W$ as above, then restrict it to the null space of activations on a
set of \emph{preservation} prompts: with $V$ the top singular vectors of
those activations, apply $\Delta W (I - V V^{\top})$
(the AlphaEdit constraint \cite{fang2025alphaedit}). Optional Gaussian layer
weighting spreads the edit around a chosen depth.

\textbf{Why it works.} If an input lies in the span of $V$, the constrained
edit has no effect on it, so behaviour on preserved inputs is guaranteed
unchanged to first order; the edit can only act on directions that harmless
inputs do not use.

\textbf{Why it exists.} It brings guarantees from knowledge editing
\cite{meng2022locating,fang2025alphaedit} to refusal removal, as an
alternative to tuning strength against a KL budget.

\textbf{Pros.} Principled protection of chosen behaviour; single pass, no
search.
\textbf{Cons.} If harmful and harmless activations share directions (likely in
newer models), the constraint also blocks the refusal edit; the rank ratio is a
new hyperparameter; dense models only in the current implementation. It is not
evaluated in this study (\S\ref{sec:m-not}).

% ---------------------------------------------------------------------
\subsection{Multi-direction and subspace methods}

\subsubsection{SVD multi-direction (Gabliteration, OBLITERATUS presets)}
\label{sec:m-svd}
\textbf{Mechanism.} Instead of the mean difference, take the top-$k$ singular
vectors of the matrix of per-prompt differences between harmful and harmless
activations, and project all $k$ out, typically with ridge regularisation and
selected layers \cite{gabliteration2024}. OBLITERATUS packages this
(``advanced'', 4 directions), whitened and iterative variants (``aggressive''),
frequency-band weighting of layers via a DCT (``spectral cascade''), and a
re-implementation of gradient-refined directions (``rdo'') in one toolkit.

\textbf{Why it works.} When refusal is spread over several correlated
directions, one mean difference captures only the dominant one; the remaining
components can re-trigger refusal. Removing a subspace closes those
alternative routes.

\textbf{Why it exists.} Single-direction ablation leaves residual refusals on
some prompt categories and some models; multi-direction methods target that
residual.

\textbf{Pros.} Removes more refusal on models where one direction is
insufficient; still closed-form and fast.
\textbf{Cons.} Every extra direction removed is extra capacity destroyed;
choosing $k$ trades refusal against damage; several presets are heuristic
combinations without a clear theoretical justification. The OBLITERATUS
v0.1.4 presets \texttt{advanced} and \texttt{spectral\_cascade}, together with
\texttt{basic}, are not evaluated in this study; they were run on
Qwen2.5-7B-Instruct and scored with this study's harness in the companion
paper \cite{young2026paper30}. Version 0.1.4 lists a \texttt{som} preset but
does not ship its direction module and does not run on Qwen3.5-4B
(Appendix~\ref{app:setup-methods}).
% src (2026-11-06, author decision: OBLITERATUS reported in paper 30 only): results/bench-4090/qwen2.5-7b__obliteratus-{basic,advanced,spectral_cascade}/push_manifest.json;
% scripts/run_external_small.sh (header NOTE: som extractor obliteratus.analysis.som_directions not shipped; v0.1.4 fails closed on every Qwen3.5); ../paper30/v3_2611/analysis/RESULTS.md

\subsubsection{Refusal cones (RDO)}
\label{sec:m-rdo}
\textbf{Mechanism.} Learn refusal directions by gradient descent on an
objective that requires ablation to remove refusal and addition to induce it,
while keeping other behaviour fixed; multiple such directions form a
polyhedral \emph{concept cone} \cite{wollschlager2025geometry}.

\textbf{Why it works / exists.} It showed that refusal is not a single
direction but a multi-dimensional cone of causally relevant, representationally
independent directions, and that directions found by optimisation outperform
diff-in-means.
\textbf{Pros.} Directly targets the causal effect rather than a statistical
proxy. \textbf{Cons.} Requires gradients and more compute. It is not evaluated
in this study (\S\ref{sec:m-not}).

\subsubsection{Self-organising-map directions (SOM)}
\label{sec:m-som}
\textbf{Mechanism.} Fit a small self-organising map (e.g.\ $4\times4$) to the
harmful activations at each layer, take $k$ map neurons as prototypes of
distinct ``kinds'' of harmful input, and use each prototype minus the harmless
mean as a refusal direction; ablate all $k$, each with its own strength
\cite{piras2025som}. The proof-of-concept Heretic integration (PR~\#196)
tunes the per-direction strengths with TPE.

\textbf{Why it works.} Different harm categories (weapons, self-harm, cyber,
\ldots) may trigger refusal through different directions; a single mean
averages them and represents none well. The SOM finds the clusters
without labels. Piras et al.~\cite{piras2025som} show diff-in-means is a special case.

\textbf{Pros.} Strong reported results (3/100 refusals at KL $\approx 0.1$ on
gpt-oss-20b, PR~\#196); principled multi-direction choice.
\textbf{Cons.} More hyperparameters (map size, $k$, per-direction strengths)
make TPE converge slowly; the PoC's top-$k$ selection counts winners over the
map's own weights instead of the data. In this study SOM is ported to
Heretic~2.0 as a modifier (``Heretic-tuned SOM'') that ranks neurons by wins
over the harmful residuals, orthonormalises the $k$ directions and replaces the
per-direction strengths with one searched decay; these and the other deviations
from Piras et al.~\cite{piras2025som} are listed in \S\ref{sec:setup-deviations}.

% ---------------------------------------------------------------------
\subsection{Optimisation-based methods}

\subsubsection{Arbitrary-Rank Ablation (ARA)}
\label{sec:m-ara}
\textbf{Mechanism.} ARA makes no assumption about the geometry of refusal
\cite{weidmann2026ara}. For each targeted linear module $W$ in a layer range
(in Heretic, the attention output projection and the MLP down-projection), it
records the module's inputs $X$ and outputs $Y = XW^{\top}$ on harmless
($g$) and harmful ($b$) prompts, then learns a low-rank update
$\Delta W = BA$ (rank 50 by default) minimising
\begin{equation}
\label{eq:ara-objective}
\mathcal{L} = w_p\,\big\|\tilde{Y}_g - Y_g\big\|^2
 + w_s\Big[\,d_k(\tilde{Y}_b, Y_g) - w_o\, d_k(\tilde{Y}_b, Y_b)\Big],
\end{equation}
where $\tilde{Y}$ is the output of $W+BA$ with each row rescaled to the norm of
the corresponding row of $W$, and $d_k(\cdot,\cdot)$ is the mean distance to the
$k$ nearest neighbours. The rescaling enters only the objective; the exported
adapter $BA$ is not renormalised (\S\ref{sec:setup-methods}).
The three terms keep harmless outputs fixed, pull harmful outputs toward
harmless ones, and push them away from where they were (``overcorrection'').
Each module is solved independently with L-BFGS; Heretic's TPE searches the
layer range, the three weights and $k$.

\textbf{Why it works.} Rather than guessing the shape of the refusal
representation, ARA directly asks each module to make harmful inputs produce
the outputs harmless inputs produce. Any refusal mechanism that a module's
output carries, whether linear, multi-directional or non-linear across
prompts, is in scope, up to the adapter's rank. The objective is close to
convex and starts at its optimum for the preservation term, so L-BFGS
converges in a few iterations.

\textbf{Why it exists.} Its author reports that directional ablation
increasingly fails on recent models, and that ARA reaches lower refusal at
lower KL on gpt-oss-20b than any directional result \cite{weidmann2026ara}.

\textbf{Pros.} No geometric assumption; fewer search parameters than
per-component directional kernels; memory similar to abliteration because
modules are solved one at a time.
\textbf{Cons.} Slower per trial; rank-50 adapters per module are larger
artefacts; the nearest-neighbour objective depends on the prompt sets, so it
may overfit them; no published analysis yet (preprint forthcoming).

% ---------------------------------------------------------------------
\subsection{Methods in other representation spaces}

\subsubsection{Jacobian-space (J-space) ablation}
\label{sec:m-jspace}
\textbf{Mechanism.} Using a Jacobian lens, map each layer's activations into a
``workspace'' where features are read out by the model's later computation;
compute the refusal direction there and pull it back to the residual stream
through the layer Jacobian, $\mathbf{d}_\ell = J_\ell^{\top}(\cdot)$; ablate it
with a forward hook rather than a weight edit
\cite{sekinal2026jspace}.

\textbf{Why it exists.} To localise the part of the refusal signal the model
actually uses downstream, rather than everything correlated with harmfulness.
\textbf{Pros.} Very targeted (small distribution shift).
\textbf{Cons.} By its authors' own measurements removes little refusal
(AdvBench refusal 0.78 versus 0.06 for diff-in-means on Qwen3.5-4B); the
reference code is hook-based, so not directly an exportable model, and is a
single-model implementation. In this study the J-space directions are instead
baked into the weights with Heretic's directional-ablation edit and search
(full row normalisation, no harmless-mean projection), which yields an
exportable model (\S\ref{sec:setup-methods}); the reference forward hook is
planned as a further condition. Because, among the candidate models,
instruction-tuned Jacobian lenses (\texttt{neuronpedia/jacobian-lens}) exist
only for the two small-pair models, J-space is run on the small pair only.

% ---------------------------------------------------------------------
\subsection{Methods not evaluated}
\label{sec:m-not}
Null-space-constrained abliteration (\S\ref{sec:m-null}) and refusal cones
(\S\ref{sec:m-rdo}) are described above for completeness but are not part of
the final design. A code audit before the main runs found that the available
null-space implementation (\texttt{jwest33/abliterator}) applies the
constraint on the output side, where it gives no preservation guarantee, and
edits square attention output projections, as in Qwen2.5-32B, OLMo-2-32B,
Olmo-3.1-32B and Mistral-Small-2409, in the wrong space; the revised
plan therefore dropped it. The only RDO implementation considered, the
\texttt{rdo} preset of OBLITERATUS, is a probe-based proxy rather than the
gradient-based optimisation of Wollschl\"ager et
al.~\cite{wollschlager2025geometry}, and it is not among the OBLITERATUS
presets evaluated in the companion paper \cite{young2026paper30}.
% src (2026-11-05 todo pass): notes/EXECUTION_PLAN_v2.md (finding 9: "jwest33 square-o_proj bug"; Design: no null-space condition; small-pair OBLITERATUS presets basic/advanced/spectral_cascade/som); notes/methods_survey_other_tools.md (jwest33 null-space constraint; OBLITERATUS presets incl. rdo); notes/session_logs/transcripts_redacted/b7a10890-24c5-4939-ba96-da845ba9da41.jsonl (code-audit report, line 4489: jwest33 "AlphaEdit-style null-space constraint, but it is applied on the output side, so it gives no preservation guarantee. Has the square o_proj bug"; "abliterate.py:2085,2157 and null_space.py:314 test the input dimension first. A square o_proj (Qwen2.5-32B, OLMo-2-32B, Olmo-3.1-32B, Mistral-Small-2409) is therefore ablated in the wrong space"; "Its rdo is a probe-proxy, not Wollschlaeger's RDO"; line 7173: null-space "dropped because its tool was buggy"). No note gives a separate, explicit reason for leaving the OBLITERATUS rdo preset out of the plan. (2026-11-06) Presets evaluated in paper 30: basic, advanced, spectral_cascade (results/bench-4090/qwen2.5-7b__obliteratus-*; ../paper30/v3_2611/analysis/RESULTS.md), no rdo.
Other diff-in-means implementations (DECCP, ErisForge, FailSpy's tool) are not
run separately; Heretic-tuned DiM represents this recipe
\cite{young2025comparative}.
Optimal-transport ablation maps the harmful activation distribution onto the
harmless one with a closed-form Gaussian transport map at one or two layers
and reports up to 11\% higher attack success than baselines on six 7--32B
models \cite{nanfack2026ot}; no public code was found. Weighted refusal
direction optimisation (WRDO) learns weights over a bank of per-prompt refusal
directions and was benchmarked against basic refusal-direction ablation, RDO,
RDO cones and Heretic on Qwen3-8B-Base and Olmo~3 7B \cite{evseev2026wrdo};
it is not evaluated. Circuit-breaker training
\cite{zou2024circuit} and extended-refusal fine-tuning
\cite{abushairah2025defense} are defences rather than removal methods and are
discussed in Section~\ref{sec:related}.

% ---------------------------------------------------------------------
% NOTE FOR AUTHORS (not for the paper): design space for new methods.
% Each method above fixes one choice on each axis; new methods = new combinations.
%  1. Where to look:   residual stream | module I/O | J-space | attention heads | experts (MoE)
%  2. What refusal is: 1 direction | subspace | cone | clusters | no assumption (learned map)
%  3. How to find it:  mean diff | SVD | SOM | gradients | direct output matching (ARA)
%  4. How to protect:  nothing | projection | norm preservation | null space | KL budget | preserve-loss
%  5. How to apply:    weight projection | low-rank adapter | hook | per-expert / router edit
%  6. How to tune:     fixed | grid | TPE | closed form
% Candidate ideas:  null-space-constrained ARA (axis 4 on ARA); per-harm-category ARA
% (SOM clusters as ARA targets); MoE router/expert-aware ARA; judge-in-the-loop objective
% instead of keywords; multi-token / sequence-level preservation loss.

